\section{Cuatro casillas y su exterior}
Quitamos las cuatro casillas del soporte y llamamos $D$ al exterior.
En la clase difusa esta supresión conserva los parámetros con margen: cada fibra pierde
$O(1/j)$, y el argumento de reequilibrio de R7, apartado 1, se aplica sin suprimir vértices.
Se mantienen ceros estructurales. Sea $x_i$ la masa que se asigna a la casilla $i$.
Cada uno de sus dos extremos especiales tiene margen exterior $t_v=1-x_i$ y los demás $t_v=1$.
Puesto $s=1/j$, escribimos los pesos de torres del exterior $p_{rc}=s t_rt_c v_{rc}$.
Denotemos por $H_{\rm ext}$ el máximo con estos márgenes, con sumando
$-p\log p+(1-p)\log(1-p)$. Definimos
\begin{align*}
 R(x)&=H_{\rm ext}(x)+\sum_v t_v\log t_v,\\
 A(x)&=\tfrac12\sum_v\log t_v-\tfrac12\sum_e\log b_2(p_e)
                     -\tfrac12\log\pdet S_{\rm ext},\\
 b_2(p)&=(1-2p)/(1-p)^2,\quad
 S_{\rm ext}=C\diag(\tau_2(p_e))C^t,\quad \tau_2(p)=p(1-p)/(1-2p).
\end{align*}
$C$ es la incidencia RC. $A(x)$ es aquí una función, distinta de la máscara $A$.
Los límites con $t_v=0$ se toman mediante la fórmula normalizada del apartado 5.
Para el conjunto $E$ de enlaces activos entre $12,34$ el óptimo se obtiene maximizando
$H_E(x)+R(x)$, donde $H_E$ es la entropía de la ley local con esas marginales
y exclusiones. El prefactor adicional es
\[
 A(x)-\tfrac12\log\det(I-V_E R''(x))+\tfrac12\log(4\pi j).
\]
Aquí $V_E$ es la covarianza local evaluada en su campo inclinado. Son las identidades
de R3, apartado 08b: al fijar las masas, la Hessiana del exterior da
$R''_{il}=\sum_{v\in i\cap l}1/t_v-u_i^t S_{\rm ext}^+u_l$.
La eliminación de las cuatro masas deja $\det(V_E^{-1}-R'')$.
En este caso el bosque local es un producto de componentes de uno o dos sitios.
Para una componente de dos sitios excluyentes, el determinante de covarianza es
$x_i x_l(1-x_i-x_l)$; para una Bernoulli vale $x_i(1-x_i)$.
Estos factores cancelan los locales de las marcas en el complemento de Schur
exactamente como en 08b. Ningún enlace RC entre dos marcas se incluye aquí.
Prolongamos las expresiones obtenidas en el interior, una vez cancelados los logaritmos
singulares. No se calcula un máximo a márgenes complejos.
